\chapter{Low Testosterone}
\index{Low Testosterone}

\section*{Definition and Overview}
Low testosterone, clinically referred to as male hypogonadism, is a biochemical and clinical syndrome characterised by a deficiency in testosterone levels, which are normally produced by the Leydig cells of the testes under the regulation of the hypothalamic-pituitary-gonadal (HPG) axis. Testosterone is the primary male sex hormone, vital for the development of male reproductive tissues, maintenance of secondary sexual characteristics, muscle mass, bone density, libido, and erythropoiesis. Hypogonadism can be primary (testicular failure) or secondary (hypothalamic or pituitary dysfunction) and can present as a congenital or adult-onset condition. In adults, low testosterone can profoundly impact quality of life and is associated with metabolic syndrome, cardiovascular disease, and osteoporosis.

\section*{Epidemiology}
Low testosterone is a prevalent condition that increases with age. Globally, it affects approximately 2\% to 6\% of middle-aged and older men. In many countries, the prevalence of biochemical androgen deficiency is estimated to affect around 10\% of men over 60, rising to over 20\% in men over 80. However, the prevalence of clinically significant hypogonadism (combining low serum levels with specific symptoms) is lower, affecting approximately 2\% of the general male population. Testosterone levels naturally decline by approximately 1\% per year starting around age 30, but rapid declines are often accelerated by co-morbidities such as obesity, type 2 diabetes, and chronic illness.

\section*{Aetiology and Pathophysiology}
Hypogonadism is categorised based on the anatomical level of the dysfunction within the HPG axis:
\begin{itemize}
    \item \textbf{Primary hypogonadism (hypergonadotropic):} Results from intrinsic testicular failure. Causes include Klinefelter syndrome (47,XXY), cryptorchidism, testicular trauma, torsion, mumps orchitis, chemotherapy, radiation, or age-related testicular decline.
    \item \textbf{Secondary hypogonadism (hypogonadotropic):} Results from hypothalamic or pituitary dysfunction, leading to inadequate secretion of gonadotropin-releasing hormone (GnRH), luteinising hormone (LH), or follicle-stimulating hormone (FSH). Causes include Kallmann syndrome, pituitary adenomas, hyperprolactinaemia, severe obesity, type 2 diabetes, chronic opioid or corticosteroid use, and haemochromatosis.
    \item \textbf{Late-onset hypogonadism:} A mixed form occurring in older men, characterised by both testicular Leydig cell depletion and blunted hypothalamic responsiveness, heavily influenced by chronic systemic disease.
    \item \textbf{Androgen receptor resistance:} Rare genetic conditions where testosterone production is normal or elevated, but target tissues cannot respond to the hormone.
\end{itemize}

\section*{Clinical Features}
The symptoms of low testosterone are diverse and can overlap with other physiological and psychological conditions. Key clinical features include:
\begin{itemize}
    \item \textbf{Sexual symptoms:} Reduced libido (sex drive), erectile dysfunction, fewer spontaneous morning erections, and difficulty achieving orgasm.
    \item \textbf{Physical changes:} Decreased muscle mass and strength, increased body fat (particularly abdominal adiposity), gynaecomastia (breast tissue enlargement), and loss of body and facial hair.
    \item \textbf{Cognitive and emotional symptoms:} Fatigue, low energy, irritability, depressed mood, difficulty concentrating, and sleep disturbances.
    \item \textbf{Skeletal changes:} Reduced bone mineral density, osteopenia, osteoporosis, and an increased risk of low-trauma fractures.
    \item \textbf{Metabolic features:} Insufficient response to exercise, insulin resistance, and dyslipidaemia.
\end{itemize}

\section*{Differential Diagnosis}
Because the symptoms of low testosterone are non-specific, several other conditions must be considered:
\begin{itemize}
    \item \textbf{Major depressive disorder:} Presents with fatigue, low mood, and reduced libido, but has normal testosterone levels.
    \item \textbf{Thyroid disorders:} Hypothyroidism can cause fatigue, weight gain, and sexual dysfunction, resembling hypogonadism.
    \item \textbf{Chronic fatigue syndrome:} Characterised by persistent, unexplained fatigue unrelated to androgen levels.
    \item \textbf{Obstructive sleep apnoea (OSA):} Causes daytime fatigue and can secondarily suppress testosterone levels due to sleep fragmentation.
    \item \textbf{Medication side effects:} Selective serotonin reuptake inhibitors (SSRIs), beta-blockers, and anti-androgens can cause sexual dysfunction and fatigue.
\end{itemize}

\section*{When to Seek Medical Care}
Men experiencing persistent sexual dysfunction, unexplained fatigue, depressed mood, or decreased physical capacity should seek medical advice.

\textbf{Contact your local emergency services for:}
\begin{itemize}
    \item \textbf{Chest pain or shortness of breath:} Sudden onset of cardiovascular symptoms, which may occur in patients with underlying metabolic syndrome.
    \item \textbf{Severe depressive symptoms:} Acute suicidal ideation or self-harm intent.
    \item \textbf{Signs of systemic crisis:} Unexplained collapse, severe headache, or sudden visual field changes (suggestive of pituitary apoplexy).
\end{itemize}
Other reasons to see a doctor include progressive loss of muscle strength, unexplained bone fractures, or difficulty conceiving.

\section*{Diagnosis and Investigations}
The diagnosis of hypogonadism must combine clear clinical features with documented biochemical deficiency:
\begin{itemize}
    \item \textbf{Total testosterone measurement:} Done via a morning blood test (between 7:00 AM and 11:00 AM) on at least two separate occasions when the patient is healthy, as levels fluctuate throughout the day and are suppressed by acute illness.
    \item \textbf{Free testosterone calculation:} Indicated if sex hormone-binding globulin (SHBG) levels are altered (e.g. in obesity or liver disease).
    \item \textbf{Luteinising hormone (LH) and FSH:} Differentiates primary (high LH/FSH) from secondary (low/normal LH/FSH) hypogonadism.
    \item \textbf{Prolactin level:} Checked to rule out hyperprolactinaemia or a pituitary prolactinoma, particularly if testosterone is severely low.
    \item \textbf{Bone densitometry (DEXA):} Indicated to assess bone mineral density in men with established hypogonadism.
    \item \textbf{Pituitary MRI:} Indicated if secondary hypogonadism is diagnosed and serum testosterone is very low (e.g. less than 5.2 nmol/L) or if visual symptoms are present.
\end{itemize}

\section*{Management}
Treatment aims to restore testosterone levels to the physiological range and relieve symptoms:
\begin{itemize}
    \item \textbf{Testosterone replacement therapy (TRT):} Administered via transdermal gels or creams, intramuscular injections (short-acting or long-acting undecanoate), or oral capsules. TRT should only be initiated when hypogonadism is pathologically confirmed.
    \item \textbf{Treating underlying causes:} Weight loss and metabolic control in obese or diabetic men, discontinuation of offending drugs, or treatment of hyperprolactinaemia.
    \item \textbf{Lifestyle measures:} Structured resistance exercise to maintain muscle mass, a balanced nutrient-dense diet, avoidance of alcohol abuse, and smoking cessation.
    \item \textbf{Monitoring during therapy:} Regular checks of hematocrit (to monitor for polycythaemia), prostate-specific antigen (PSA), digital rectal examination (DRE), and serum testosterone levels to ensure therapeutic target ranges are met.
\end{itemize}

\section*{Complications}
Complications of low testosterone and its treatment include:
\begin{itemize}
    \item \textbf{Cardiovascular disease:} Association between low testosterone and increased cardiovascular risk, though TRT safety must be evaluated individually.
    \item \textbf{Polycythaemia (erythrocytosis):} Excessive red blood cell production, a common side effect of TRT that increases blood viscosity.
    \item \textbf{Infertility:} Exogenous testosterone suppresses spermatogenesis; alternative treatments (e.g. hCG or clomiphene) must be used if fertility is desired.
    \item \textbf{Prostate events:} While TRT does not cause prostate cancer, it can stimulate the growth of existing hormone-sensitive prostate lesions.
    \item \textbf{Sleep apnoea exacerbation:} TRT can worsen obstructive sleep apnoea in susceptible individuals.
\end{itemize}

\section*{Prognosis and Outcomes}
The prognosis for men with low testosterone is highly favourable when treated appropriately. TRT generally leads to significant improvements in libido, sexual function, mood, energy levels, and body composition (increased muscle and reduced fat) within 3 to 6 months. Bone mineral density improves progressively over 12 to 24 months. For secondary hypogonadism caused by lifestyle factors (such as obesity), substantial weight loss and metabolic optimisation can restore testosterone levels without the need for lifelong hormone replacement.

\section*{Prevention and Self-Care}
While primary hypogonadism cannot be prevented, late-onset hypogonadism can be mitigated or managed at home:
\begin{itemize}
    \item \textbf{Maintain a healthy weight:} Obesity increases the conversion of testosterone to oestrogen in adipose tissue.
    \item \textbf{Regular physical activity:} Incorporate strength and resistance training, which naturally supports muscle mass and metabolic health.
    \item \textbf{Manage chronic diseases:} Keep blood glucose and blood pressure well controlled to protect Leydig cell function.
    \item \textbf{Prioritise sleep hygiene:} Testosterone production peaks during REM sleep; chronic sleep deprivation lowers hormone levels.
    \item \textbf{Limit alcohol consumption:} Excessive alcohol intake acts as a direct testicular toxin and disrupts pituitary function.
\end{itemize}

\section*{Sources and Further Reading}
The following reputable organisations provide further information on testosterone deficiency:
\begin{itemize}
    \item Healthy Male (Andrology Australia). \textit{Testosterone deficiency and hypogonadism resources.} https://www.healthymale.org.au/
    \item Endocrine Society of Australia (ESA). \textit{Clinical guidelines for male hypogonadism.} https://www.endocrinesociety.org.au/
    \item Better Health Channel. \textit{Testosterone deficiency in men.} https://www.betterhealth.vic.gov.au/
    \item Mayo Clinic. \textit{Testosterone therapy: key clinical considerations.} https://www.mayoclinic.org/
    \item NHS. \textit{Male menopause and testosterone replacement.} https://www.nhs.uk/
\end{itemize}
